\section*{第 \ref{ch:b3:microbiomes} 章——微生物组与共生}

\begin{solution}{exo:b3:microbiomes:1}
\omterm{def:b3:microbiomes:symbiosis}{互利共生}：双方都受益——豆科植物与\omterm{def:b3:microbiomes:nodule}{根瘤菌}，珊瑚与藻类。
\omterm{def:b3:microbiomes:symbiosis}{偏利共生}：一方受益，另一方不受影响——大多数肠道细菌
以宿主无法消化的物质为食。寄生：一方以
另一方为代价而受益——\emph{Wolbachia} 杀死雄性胚胎。同一对
可以转变：肠道\omterm{def:b3:microbiomes:symbiosis}{共生}菌穿过肠壁后成为病原体，
停止固氮的\omterm{def:b3:microbiomes:nodule}{根瘤菌}成为\omterm{def:b3:microbiomes:nodule}{根瘤}的寄生生物，而
热胁迫下的藻类会伤害容纳它的珊瑚。
\end{solution}

\begin{solution}{exo:b3:microbiomes:2}
$H = -(0.5\ln 0.5 + 0.3\ln 0.3 + 0.1\ln 0.1 + 2\times 0.05\ln 0.05) =
0.347 + 0.361 + 0.230 + 0.300 = 1.24$；$e^{H} = 3.4$ 有效
物种。$D = 0.25 + 0.09 + 0.01 + 0.0025 + 0.0025 = 0.355$；$1/D =
2.8$。
\end{solution}

\begin{solution}{exo:b3:microbiomes:3}
16S 调查报告了有哪些微生物存在，大致到属水平，以及它们的
相对读段数，但不能说明它们携带的基因，
漏掉真菌和\omterm{def:b3:virology:virus}{病毒}（需用其他\omterm{def:b3:genetic-engineering:pcr}{引物}），并受拷贝数
和\omterm{def:b3:genetic-engineering:pcr}{引物}的影响。鸟枪法宏基因组测序报告基因、功能，以及
在足够深度下得到的完整基因组和菌株；每个样本成本更高，
组织样本中易被宿主 DNA 淹没，仍无法判断哪些基因
正在表达，或哪些细胞是活的。
\end{solution}

\begin{solution}{exo:b3:microbiomes:4}
将纤维发酵为\omterm{def:b3:microbiomes:gut}{短链脂肪酸}，为结肠和
宿主提供营养；合成维生素~K 和 B 族维生素；转化胆汁
酸和药物；提供\omterm{def:b3:microbiomes:gut}{定植抗性}，抵御病原体；训练
免疫系统。\omterm{def:b3:microbiomes:gut}{菌群失调}的例子是：\emph{Clostridioides difficile} 结肠炎
（发生于\omterm{def:b3:bacteriology:antibiotics}{抗生素}之后；另有关联：炎症性肠病、
肥胖）。
\end{solution}

\begin{solution}{exo:b3:microbiomes:5}
$\binom{20}{5} = \num{15504}$。物种数为 $10$ 时：$1 - \binom{10}{5}/
\num{15504} = 1 - 252/\num{15504} = 0.984$；物种数为 $6$ 时：$1 - 2002/\num{15504}
= 0.871$；物种数为 $3$ 时：$1 - 6188/\num{15504} = 0.601$；物种数为 $1$ 时：$1 -
\num{11628}/\num{15504} = 0.25$。$E[S_{5}] = 2.7$ 个物种。覆盖度：
一个单例，$1 - 1/20 = 0.95$。
\end{solution}

\begin{solution}{exo:b3:microbiomes:6}
$400\times 10^{11} = 4\times 10^{13}$ 个细菌，约 \qty{40}{g}。基因：
$1000$ 个物种 $\times$ \num{4000}，扣除重叠后——数量级约为
$3\times 10^{6}$ 个不同基因，是人类的
\num{20000} 个基因的一百多倍。
\end{solution}

\begin{solution}{exo:b3:microbiomes:7}
\omterm{def:b3:microbiomes:nodule}{固氮酶}会被氧气破坏，但类菌体必须进行呼吸，
以便每固定一个 N$_{2}$ 产生十六个 ATP。\omterm{def:b3:microbiomes:nodule}{根瘤}皮层形成一道扩散
屏障，限制氧气进入，而\omterm{def:b3:microbiomes:nodule}{豆血红蛋白}以毫摩尔浓度存在，
结合进入的氧气，使游离浓度
维持在 \qty{10}{nM} 附近，同时流向类菌体的通量
很高。粉红色为氧合\omterm{def:b3:microbiomes:nodule}{豆血红蛋白}；绿色\omterm{def:b3:microbiomes:nodule}{根瘤}已降解了其
\omterm{def:b3:microbiomes:nodule}{豆血红蛋白}——已\omterm{def:b3:cell-cycle-apoptosis:other}{衰老}，不再固氮。
\end{solution}

\begin{solution}{exo:b3:microbiomes:8}
两者均失效：既无\omterm{def:b3:microbiomes:nodule}{根瘤}，也无\omterm{def:b3:microbiomes:mycorrhiza}{丛枝菌根}，因为
钙峰是\omterm{def:b3:microbiomes:mycorrhiza}{共同共生途径}共有的信号。
因此该途径早于\omterm{def:b3:microbiomes:nodule}{根瘤}的形成——它最初为
四亿年前的真菌\omterm{def:b3:microbiomes:symbiosis}{共生}而建，后来被豆科植物
招募，用于与\omterm{def:b3:microbiomes:nodule}{根瘤菌}\omterm{def:b3:microbiomes:symbiosis}{共生}。
\end{solution}

\begin{solution}{exo:b3:microbiomes:9}
基尔斯将部分大豆\omterm{def:b3:microbiomes:nodule}{根瘤}置于氩气与
氧气的气体环境中，\omterm{def:b3:microbiomes:nodule}{根瘤菌}无法固氮，而该植株
其余\omterm{def:b3:microbiomes:nodule}{根瘤}正常固氮；植株切断了不固氮\omterm{def:b3:microbiomes:nodule}{根瘤}的供氧，
这些\omterm{def:b3:microbiomes:nodule}{根瘤}中的\omterm{def:b3:microbiomes:nodule}{根瘤菌}繁殖量仅约为
一半。若无此区分，跳过昂贵固氮过程的
\omterm{def:b3:microbiomes:nodule}{根瘤菌}会比固氮者繁殖更快，并在
种群中扩散，最终植株将容纳不给予任何回报的细菌，
\omterm{def:b3:microbiomes:symbiosis}{互利共生}将随之瓦解。
\end{solution}

\begin{solution}{exo:b3:microbiomes:10}
$\binom{N-N_{i}}{n}/\binom{N}{n} = \prod_{j=0}^{n-1}(N - N_{i} - j)/(N -
j)$，每个因子均小于 $1$；由 $n$ 增至 $n + 1$ 时，多乘上一个
小于 $1$ 的因子，于是漏检概率下降，
存在概率随 $n$ 增大而上升；当 $n = N$ 时，分子 $\binom{N -
N_{i}}{N}$ 为零，每个物种都存在，结果为 $S$。对于 $N_{i},
n \ll N$ 时，每个因子约为 $1 - N_{i}/N$，乘积约为 $(1 -
N_{i}/N)^{n} \approx e^{-nN_{i}/N} \approx (1 - n/N)^{N_{i}}$——此即
有放回抽样近似。
\end{solution}

\begin{solution}{exo:b3:microbiomes:11}
雄性对于只通过卵传播的\omterm{def:b3:microbiomes:symbiosis}{共生}体而言是死胡同，因此任何
使雄性转变为雌性、杀死雄性以利于其姐妹，
或使未感染的雌性成为累赘，都会提高\omterm{def:b3:microbiomes:symbiosis}{共生}体的
传递率。胞质不亲和：感染雄性的精子经过
改造，使受精卵死亡，除非卵子携带相同的
\emph{Wolbachia}，可拯救受精卵。未感染的雌性因此损失
与感染雄性所孕育的后代，而感染雌性
则不损失任何后代；优势随感染频率增大而增强，因此超过
一定阈值后，感染会扩散至固定，即使感染雌性
产卵略少——这一代价被
未感染雌性后代被毁的比例所抵消。
\end{solution}

\begin{solution}{exo:b3:microbiomes:12}
内\omterm{def:b3:microbiomes:symbiosis}{共生}体每经历一代宿主，都要经过少量细胞的瓶颈，
且从不与外来者重组：漂变固定
轻度有害的突变（包括缺失）被固定，
且没有替代基因的来源（穆勒棘轮）。其产物由宿主提供的基因，
或服务于自由生活的基因——\omterm{def:b3:virology:virus}{包膜}、运动性、
调控、多数修复相关基因——不会被察觉缺失，最先丢失；
翻译机器，以及制造宿主所需物质的基因
（\emph{Buchnera} 中的必需氨基酸）保留最久。自由生活的
近缘类群种群庞大，有重组和基因流入，
自然选择维持着它们的基因组。逆转需要外部的
新基因——在孤立时不可能获得——因此宿主转而
以新的\omterm{def:b3:microbiomes:symbiosis}{共生}体替换已退化者，正如若干昆虫谱系所做的那样。
\end{solution}

\begin{solution}{pb:b3:microbiomes:1}
\textbf{1.} $400\times 10^{11} = 4\times 10^{13}$ 个细菌，\qty{40}{g}。
\textbf{2.} 含 \num{4000} 个基因的基因组约为 \qty{4}{Mb}，\qty{4.4}{fg}：
细菌 DNA 质量为 $4\times 10^{13}\times 4.4\times 10^{-15} = \qty{0.18}{g}$，
对比 $3\times 10^{13}\times 6.4\times 10^{-12} = \qty{0.19}{g}$ 的
人体 DNA——大致相当。
\textbf{3.} $1000\times 0.7\times 4000 + 0.3\times 4000 \approx 2.8\times
10^{6}$ 个不同基因，约为人类基因组的 $140$ 倍。
\textbf{4.} $30\times 10 = \qty{300}{mmol}$; $\times 0.9 = \qty{270}{kJ}$;
占饮食的 \qty{3}{\%}。
\textbf{5.} 否——\qty{3}{\%} 不足以解释 \qty{30}{\%}。无菌
小鼠的吸收效率也较低，肠道激素、脂肪
储存与产热均有改变，食欲也不同：\omterm{def:b3:microbiomes:symbiosis}{微生物组}
作用于宿主的生理功能，而不仅是燃料的供应者。
\textbf{6.} 分裂与排出相平衡：种群稳定，每个细胞
每日被替换。经 \qty{99.9}{\%} 杀灭后尚余
$4\times 10^{10}$ 个，十次倍增可在约十天内恢复数量，但存活者与
生长最快者占主导，组成发生改变（即\omterm{def:b3:microbiomes:gut}{菌群失调}），
在此间隙，如 \emph{C.
difficile} 一类入侵者即可定居。
\textbf{7.} 次要物种合计占 $0.4$，由 $177$ 个物种分担，每个 $p = 0.00226$。$H
= -(0.3\ln 0.3 + 0.2\ln 0.2 + 0.1\ln 0.1 + 0.4\ln 0.00226) = 0.361 +
0.322 + 0.230 + 2.437 = 3.35$; $e^{H} \approx 28$.
\textbf{8.} $D = 0.09 + 0.04 + 0.01 + 177\times (0.00226)^{2} = 0.141$;
$1/D = 7.1$。\omterm{def:b3:microbiomes:diversity}{辛普森指数}对丰度取平方，因此三个
优势物种决定其值，稀有物种几乎不计；\omterm{def:b3:microbiomes:diversity}{香农
指数}则更重视稀有物种。
\textbf{9.} $1 - 60/2000 = 0.97$：约每千条读段中有 $30$ 条
属于尚未见到的物种。
\textbf{10.} 丰富度随样本量增加而升高，稀有物种不断
出现；\num{10000} 条读段可覆盖 \num{2000} 条读段所漏掉的物种。
将大样本\omterm{def:b3:microbiomes:diversity}{稀疏化}至 \num{2000} 条读段（依定理，或
通过子抽样）——若为同一群落，应得约 $180$ 种，
并在同等深度下比较。
\textbf{11.} $N_{i} = 1$: $1 - \binom{1999}{1000}/\binom{2000}{1000} =
1 - (2000 - 1000)/2000 = 0.5$。当 $N_{i} = 5$ 时，约为 $1 - 0.5^{5} = 0.97$。
\textbf{12.} 有效物种数 $e^{2.0} = 7.4$；$D \ge 0.5^{2} = 0.25$，
约为 $0.27$：即一个优势物种与几十个次要物种并存的群落，
也就是崩溃的低多样性肠道，
\emph{C.~difficile} 正在其中站稳脚跟。
\textbf{13.} 每公顷糖类 $150\times 8 = \qty{1200}{kg}$，
约占光合固定量 \qty{12}{t} 的 \qty{10}{\%}。
\textbf{14.} $150\,000/28 = 5360$ 摩尔 N$_{2}$；$16\times 5360 =
\num{85700}$ 摩尔 ATP。
\textbf{15.} $\num{85700}/30 = 2860$ 摩尔葡萄糖，\qty{514}{kg}——
约占所付糖量的 \qty{43}{\%}。其余
用于构建与维持\omterm{def:b3:microbiomes:nodule}{根瘤}和类菌体，为每个 N$_{2}$ 提供八个
电子作为还原剂，并同化氨。
\textbf{16.} $1200\times 16 = \qty{19200}{MJ}$，对应 \qty{150}{kg} 的 N：
\qty{128}{MJ/kg}，几乎是工厂 \qty{35}{MJ/kg} 的四倍——但
所耗为阳光，在田间进行，无需燃料。
\textbf{17.} 固氮者 $0.9\times 1 = 0.9$；\omterm{prop:b3:microbiomes:sanctions}{欺骗者} $0.1\times 0.5 =
0.05$；\omterm{prop:b3:microbiomes:sanctions}{欺骗者}占 $0.05/0.95 = 5.3\,\%$——一季内减半。
\textbf{18.} \omterm{prop:b3:microbiomes:sanctions}{欺骗者} $0.1\times 1.2^{10} = 0.62$，对比固氮者 $0.9$：
十个季节后占 \qty{41}{\%}，并趋向占据优势。\omterm{prop:b3:microbiomes:sanctions}{制裁}使
选择转而不利于\omterm{prop:b3:microbiomes:sanctions}{欺骗者}；无此\omterm{prop:b3:microbiomes:sanctions}{制裁}，\omterm{def:b3:microbiomes:symbiosis}{互利共生}便会退化。
\textbf{19.} $10^{6}\times \qty{20}{pg} = \qty{20}{\micro g}$ 的碳，
每 cm$^{2}$ 每天；\qty{18}{\micro g} 输给珊瑚。
\textbf{20.} 呼吸作用消耗 \qty{15}{\micro g}：盈余 \qty{3}{\micro g}
每 cm$^{2}$ 每天，用于生长、钙化的有机基质，
黏液与配子。
\textbf{21.} 若藻类仅占 \qty{5}{\%}，则每天获得 \qty{0.9}{\micro g}，
而呼吸消耗 \qty{15}{\micro g}，每天亏缺 \qty{14}{\micro g}；
\qty{300}{\micro g} 的储备可维持约三周。
\textbf{22.} 在 $+2$ 下持续四周更糟：藻类
光系统的\omterm{def:b3:dna-repair:lesions}{损伤}随温度陡升，而非线性增加。该指标
仍可作为预警，因为白化综合了累积
胁迫，且两条曲线之差小于阈值的不确定度；
它是一个经验指数，依据观察到的白化情况校准。
\textbf{23.} 受胁迫的珊瑚可转向已以低丰度存在的更耐热藻类
物种，或在白化后从海水中获取它们；
耐热藻类提供的碳较少，因此珊瑚生长更慢，
所获得的耐受性仅为一至两度——
远小于预计的升温幅度，
而升温速度超过了珊瑚自身的适应能力。
\textbf{24.} $10^{10}$ cm$^{2}$ $\times$ \qty{20}{\micro g} $=
\qty{200}{kg}$ 的碳，每天，约合每年 \qty{73}{t}。
\textbf{25.} 约 \qty{3}{\%} 的饮食来自\omterm{def:b3:microbiomes:symbiosis}{微生物组}的
发酵；测序样本的覆盖度为 \qty{97}{\%}；
\qty{10}{\%} 的光合作用产物用于支付豆类作物的氮。
\end{solution}
